% theme: position_vectors_intersections_and_area_ratios
% 叩き台。交点の位置ベクトル、係数付きベクトル条件、面積比を扱う。
\MajorQuestionLesson{位置ベクトルと交点・面積比}
\SetRowSpace{2.6mm}

% figure: 三角形 ABC、P は AB 上、Q は AC 上、R=CP∩BQ
\TypeQuestion{$\Tri{ABC}$ で $\overrightarrow{\mathrm{AB}}=\vec{a}$、$\overrightarrow{\mathrm{AC}}=\vec{b}$ とする。辺 $\mathrm{AB}$ を $2:1$ に内分する点を $\mathrm{P}$、辺 $\mathrm{AC}$ を $1:2$ に内分する点を $\mathrm{Q}$ とし、線分 $\mathrm{CP}$ と線分 $\mathrm{BQ}$ の交点を $\mathrm{R}$ とする。また、$s,t$ を実数とする。}{intersection_of_cevians_basic}
\QQ{$\overrightarrow{\mathrm{AP}},\overrightarrow{\mathrm{AQ}}$ を $\vec{a},\vec{b}$ を用いて表しなさい。}{}
\QQ{点 $\mathrm{R}$ が線分 $\mathrm{CP}$ を $s:1-s$ に内分するとき、$\overrightarrow{\mathrm{AR}}$ を $s,\vec{a},\vec{b}$ を用いて表しなさい。}{}
\QQ{点 $\mathrm{R}$ が線分 $\mathrm{BQ}$ を $t:1-t$ に内分するとき、$\overrightarrow{\mathrm{AR}}$ を $t,\vec{a},\vec{b}$ を用いて表しなさい。}{}
\QQ{実数 $s,t$ と $\overrightarrow{\mathrm{AR}}$ を求めなさい。}{}
\QQ{$\Lseg{CR}:\Lseg{RP}$ と $\Lseg{BR}:\Lseg{RQ}$ をそれぞれ求めなさい。}{}

% figure: Q は AC の C 側の延長上、R=AP∩BQ
\TypeQuestion{$\Tri{ABC}$ で $\overrightarrow{\mathrm{AB}}=\vec{a}$、$\overrightarrow{\mathrm{AC}}=\vec{b}$ とする。辺 $\mathrm{BC}$ を $1:2$ に内分する点を $\mathrm{P}$、線分 $\mathrm{AC}$ を $3:1$ に外分する点を $\mathrm{Q}$ とし、直線 $\mathrm{AP}$ と直線 $\mathrm{BQ}$ の交点を $\mathrm{R}$ とする。}{intersection_with_external_division_point}
\QQ{$\overrightarrow{\mathrm{AP}},\overrightarrow{\mathrm{AQ}}$ を $\vec{a},\vec{b}$ を用いて表しなさい。}{}
\QQ{$\overrightarrow{\mathrm{AR}}=s\overrightarrow{\mathrm{AP}}$ とするとき、$\overrightarrow{\mathrm{AR}}$ を実数 $s$ と $\vec{a},\vec{b}$ を用いて表しなさい。}{}
\QQ{$\overrightarrow{\mathrm{BR}}=t\overrightarrow{\mathrm{BQ}}$ とするとき、$\overrightarrow{\mathrm{AR}}$ を実数 $t$ と $\vec{a},\vec{b}$ を用いて表しなさい。}{}
\QQ{実数 $s,t$ と $\overrightarrow{\mathrm{AR}}$ を求めなさい。}{}
\QQ{$\Lseg{AP}:\Lseg{PR}$ と $\Lseg{BR}:\Lseg{RQ}$ を求めなさい。}{}

% figure: 平行四辺形 ABCD、E は BC 上、F は CD 上、P=AE∩BF
\TypeQuestion{平行四辺形 $\mathrm{ABCD}$ で、$\overrightarrow{\mathrm{AB}}=\vec{a}$、$\overrightarrow{\mathrm{AD}}=\vec{b}$ とする。辺 $\mathrm{BC}$ を $2:1$ に内分する点を $\mathrm{E}$、辺 $\mathrm{CD}$ を $1:2$ に内分する点を $\mathrm{F}$ とし、線分 $\mathrm{AE}$ と線分 $\mathrm{BF}$ の交点を $\mathrm{P}$ とする。}{intersection_in_parallelogram}
\QQ{$\overrightarrow{\mathrm{AE}},\overrightarrow{\mathrm{BF}}$ を $\vec{a},\vec{b}$ を用いて表しなさい。}{}
\QQ{$\overrightarrow{\mathrm{AP}}$ を $\vec{a},\vec{b}$ を用いて表しなさい。}{}
\QQ{$\Lseg{AP}:\Lseg{PE}$ と $\Lseg{BP}:\Lseg{PF}$ を求めなさい。}{}

% figure: 平行四辺形 ABCD、E は BC 上、P=DE∩AC
\TypeQuestion{平行四辺形 $\mathrm{ABCD}$ で、$\overrightarrow{\mathrm{AB}}=\vec{a}$、$\overrightarrow{\mathrm{AD}}=\vec{b}$ とする。辺 $\mathrm{BC}$ を $1:2$ に内分する点を $\mathrm{E}$ とし、線分 $\mathrm{DE}$ と対角線 $\mathrm{AC}$ の交点を $\mathrm{P}$ とする。}{intersection_with_diagonal_in_parallelogram}
\QQ{$\overrightarrow{\mathrm{AE}}$ を $\vec{a},\vec{b}$ を用いて表しなさい。}{}
\QQ{$\overrightarrow{\mathrm{AP}}=s\overrightarrow{\mathrm{AC}}$ とするとき、$\overrightarrow{\mathrm{AP}}$ を実数 $s$ と $\vec{a},\vec{b}$ を用いて表しなさい。}{}
\QQ{$\overrightarrow{\mathrm{DP}}=t\overrightarrow{\mathrm{DE}}$ とするとき、$\overrightarrow{\mathrm{AP}}$ を実数 $t$ と $\vec{a},\vec{b}$ を用いて表しなさい。}{}
\QQ{実数 $s,t$ と $\overrightarrow{\mathrm{AP}}$ を求めなさい。}{}
\QQ{$\Lseg{AP}:\Lseg{PC}$ と $\Lseg{DP}:\Lseg{PE}$ を求めなさい。}{}

\LessonPageBreak
\SetRowSpace{4mm}

% 図は不要。生徒が必要なら書く前提。
\TypeQuestion{点 $\mathrm{P}$ は $\overrightarrow{\mathrm{PA}}+2\overrightarrow{\mathrm{PB}}+3\overrightarrow{\mathrm{PC}}=\vec{0}$ を満たす。点 $\mathrm{A},\mathrm{B},\mathrm{C}$ の位置ベクトルをそれぞれ $\vec{a},\vec{b},\vec{c}$ とする。}{weighted_point_inside_triangle}
\QQ{点 $\mathrm{P}$ の位置ベクトル $\vec{p}$ を $\vec{a},\vec{b},\vec{c}$ を用いて表しなさい。}{}
\QQ{直線 $\mathrm{AP}$ と辺 $\mathrm{BC}$ の交点を $\mathrm{D}$ とするとき、$\Lseg{BD}:\Lseg{DC}$ を求めなさい。}{}
\QQ{面積比 $\Tri{PBC}:\Tri{PCA}:\Tri{PAB}$ を求めなさい。}{}

% 図は不要。生徒に書かせて。
\TypeQuestion{点 $\mathrm{P}$ は $2\overrightarrow{\mathrm{PA}}-\overrightarrow{\mathrm{PB}}+3\overrightarrow{\mathrm{PC}}=\vec{0}$ を満たす。点 $\mathrm{A},\mathrm{B},\mathrm{C}$ の位置ベクトルをそれぞれ $\vec{a},\vec{b},\vec{c}$ とする。}{weighted_point_outside_triangle}
\QQ{点 $\mathrm{P}$ の位置ベクトル $\vec{p}$ を $\vec{a},\vec{b},\vec{c}$ を用いて表しなさい。}{}
\QQ{直線 $\mathrm{AP}$ と直線 $\mathrm{BC}$ の交点を $\mathrm{D}$ とするとき、点 $\mathrm{D}$ は線分 $\mathrm{BC}$ を何対何に外分するか求めなさい。}{}
\QQ{点 $\mathrm{P}$ の位置を図示しなさい。}{}
\QQ{面積比 $\Tri{PBC}:\Tri{PCA}:\Tri{PAB}$ を求めなさい。}{}

% 図は不要。生徒が必要なら書く前提。
\TypeQuestion{$\mathrm{A}(0,0),\mathrm{B}(6,0),\mathrm{C}(0,4)$ とし、点 $\mathrm{P}$ は $2\overrightarrow{\mathrm{PA}}+\overrightarrow{\mathrm{PB}}+3\overrightarrow{\mathrm{PC}}=\vec{0}$ を満たす。}{weighted_point_area_with_components_inside}
\QQ{点 $\mathrm{P}$ の座標を求めなさい。}{}
\QQ{直線 $\mathrm{AP}$ と辺 $\mathrm{BC}$ の交点を $\mathrm{D}$ とするとき、$\Lseg{BD}:\Lseg{DC}$ を求めなさい。}{}
\QQ{$\Tri{PBC},\Tri{PCA},\Tri{PAB}$ の面積を、成分を用いて求めなさい。}{}

% 図は不要。生徒に書かせて。
\TypeQuestion{$\mathrm{A}(0,0),\mathrm{B}(4,0),\mathrm{C}(0,6)$ とし、点 $\mathrm{P}$ は $-\overrightarrow{\mathrm{PA}}+3\overrightarrow{\mathrm{PB}}+2\overrightarrow{\mathrm{PC}}=\vec{0}$ を満たす。}{weighted_point_area_with_components}
\QQ{点 $\mathrm{P}$ の座標を求めなさい。}{}
\QQ{直線 $\mathrm{AP}$ と辺 $\mathrm{BC}$ の交点を $\mathrm{D}$ とするとき、$\Lseg{BD}:\Lseg{DC}$ を求めなさい。}{}
\QQ{点 $\mathrm{P}$ の位置を座標平面上に図示しなさい。}{}
\QQ{$\Tri{PBC},\Tri{PCA},\Tri{PAB}$ の面積を、成分を用いて求めなさい。}{}
